\section{Lộ trình phát triển sản phẩm trong các giai đoạn tiếp theo}
Ở giai đoạn hiện tại, Mutux vẫn đang trong quá trình hình thành ý tưởng và xây dựng phiên bản đầu tiên của sản phẩm. Mục tiêu quan trọng nhất của nhóm không phải là phát triển một hệ thống với đầy đủ chức năng ngay từ đầu, mà là nhanh chóng xây dựng một Minimum Viable Product (MVP) để kiểm chứng tính khả thi của ý tưởng, thu thập phản hồi từ người dùng và từng bước hoàn thiện mô hình kinh doanh.

Lộ trình phát triển của Mutux được định hướng theo bốn giai đoạn như sau.

\subsection{Giai đoạn 1: Hoàn thiện MVP và kiểm chứng ý tưởng}
Trong phạm vi môn học, sản phẩm sẽ được phát triển theo hướng hoàn thiện giao diện người dùng trước, sau đó xây dựng đầy đủ các chức năng Front-end và Back-end để tạo thành một MVP hoàn chỉnh. Phiên bản này tập trung vào những chức năng cốt lõi như tìm kiếm thiết bị, đăng tải gaming gear, đặt thuê, quản lý đơn thuê và các chức năng quản trị cơ bản.

Song song với việc phát triển sản phẩm, nhóm cũng hoàn thiện Business Plan nhằm xác định rõ giá trị cốt lõi, mô hình doanh thu, khách hàng mục tiêu và chiến lược phát triển của Mutux. Đây là nền tảng quan trọng để dự án có thể tiếp tục phát triển sau khi kết thúc môn học.

\subsection{Giai đoạn 2: Thử nghiệm thị trường}
Sau khi MVP hoàn thiện, Mutux sẽ triển khai thử nghiệm với một số cửa hàng gaming gear hoặc cộng đồng game thủ có quy mô nhỏ. Mục tiêu của giai đoạn này là đánh giá khả năng sử dụng của hệ thống, mức độ phù hợp của các chức năng hiện có và xác định những khó khăn trong quá trình vận hành thực tế.

Các phản hồi thu thập được sẽ được sử dụng để cải tiến giao diện, tối ưu quy trình thuê thiết bị, bổ sung những tính năng còn thiếu và điều chỉnh mô hình kinh doanh trước khi mở rộng quy mô.

Ở giai đoạn này, số lượng người dùng chưa phải yếu tố quan trọng nhất; thay vào đó, nhóm tập trung xây dựng một sản phẩm hoạt động ổn định và đáp ứng đúng nhu cầu của thị trường.

\subsection{Giai đoạn 3: Mở rộng quy mô vận hành}
Sau khi sản phẩm đã được kiểm chứng và có những người dùng đầu tiên, Mutux sẽ bắt đầu mở rộng quy mô hoạt động. Khi số lượng người dùng tăng lên, hệ thống cần được nâng cấp về hạ tầng kỹ thuật, khả năng chịu tải và độ ổn định.

Đồng thời, nhóm sẽ từng bước chuẩn hóa quy trình vận hành như xác minh danh tính người dùng, quản lý thiết bị, xử lý tranh chấp, thanh toán và chăm sóc khách hàng. Đây cũng là thời điểm cần xây dựng đầy đủ các quy định pháp lý giữa nền tảng, người thuê và chủ sở hữu thiết bị nhằm giảm thiểu rủi ro trong quá trình kinh doanh.

Nếu lượng người dùng đạt khoảng vài nghìn tài khoản hoạt động thường xuyên, Mutux sẽ xem xét đăng ký hoạt động dưới hình thức hộ kinh doanh hoặc công ty trách nhiệm hữu hạn để phù hợp với quy mô vận hành.

\subsection{Giai đoạn 4: Phát triển doanh nghiệp}
Khi thị trường đã ổn định và số lượng giao dịch tăng trưởng bền vững, Mutux sẽ mở rộng mạng lưới đối tác bao gồm các cửa hàng gaming gear lớn, nhà phân phối và các đơn vị cung cấp dịch vụ thanh toán, vận chuyển hoặc tài chính.

Ở giai đoạn này, doanh nghiệp sẽ cân nhắc giữa hai chiến lược phát triển. Phương án thứ nhất là tiếp tục vận hành theo mô hình bootstrapping nhằm duy trì quyền kiểm soát doanh nghiệp. Phương án thứ hai là kêu gọi vốn từ các nhà đầu tư nếu cần nguồn lực để mở rộng thị trường, phát triển ứng dụng di động, đầu tư hạ tầng công nghệ hoặc mở rộng sang các phân khúc khách hàng mới.

Mục tiêu dài hạn của Mutux là trở thành một marketplace chuyên biệt dành cho hoạt động thuê và trải nghiệm gaming gear tại Việt Nam.

\subsection{Định hướng phát triển tổ chức và nhân sự}
Đội ngũ sáng lập Mutux hiện gồm năm thành viên đều có chuyên môn về Công nghệ thông tin. Đây là lợi thế trong việc phát triển sản phẩm và triển khai hệ thống phần mềm. Tuy nhiên, để vận hành một doanh nghiệp thực tế, nhóm vẫn còn thiếu kinh nghiệm ở các lĩnh vực như tài chính, marketing, phát triển kinh doanh và pháp lý.

Trong giai đoạn đầu, nhóm sẽ tận dụng mạng lưới cố vấn từ giảng viên, các chuyên gia và đối tác bên ngoài để hỗ trợ ra quyết định trong những lĩnh vực chưa có nhiều kinh nghiệm. Khi quy mô hoạt động mở rộng, Mutux sẽ ưu tiên tuyển dụng các vị trí Business Development, Marketing, Finance và Legal nhằm chuyên môn hóa quá trình vận hành, mở rộng thị trường và giảm thiểu các rủi ro pháp lý.

\subsection{Gọi vốn}
Trong giai đoạn đầu, Mutux định hướng phát triển theo mô hình bootstrapping, sử dụng nguồn lực của các thành viên để xây dựng và kiểm chứng sản phẩm. Cách tiếp cận này giúp nhóm duy trì quyền chủ động trong quá trình phát triển và hạn chế áp lực từ các nhà đầu tư khi mô hình kinh doanh chưa được kiểm chứng.

Nếu sản phẩm đạt được lượng người dùng và doanh thu ổn định, nhóm sẽ cân nhắc kêu gọi vốn nhằm mở rộng hạ tầng kỹ thuật, phát triển ứng dụng di động, mở rộng mạng lưới đối tác và nâng cao năng lực vận hành.

Về cơ cấu sở hữu, trong giai đoạn đầu quyền lợi của các thành viên sẽ được xác định thông qua thỏa thuận nội bộ và cam kết bằng văn bản, dựa trên mức độ đóng góp và vai trò của từng người. Khi doanh nghiệp phát triển thành công ty cổ phần hoặc có sự tham gia của nhà đầu tư, cơ cấu cổ phần sẽ được điều chỉnh theo quy định của pháp luật, đồng thời đảm bảo cân bằng lợi ích giữa các nhà sáng lập, nhà đầu tư và định hướng phát triển lâu dài của Mutux.